\chapter{推荐评价与效果识别}
\label{chap:rs-evaluation}

一个新模型把通勤背包排得更靠前，另一个方案减少重复展示，第三个方案让广告预算更晚耗尽。怎样判断它们是否有效？把三个方案都放进同一张命中率表，会丢掉它们实际改变的东西：模型可以在固定判断集合上排序，展示方案会改变用户看见什么，预算策略还会改变后续机会。离线成绩提高之后，仍需解释它估计的是哪一种质量，以及这些数据为什么能够支持部署后的效果判断。

前面各章已经定义候选、分数、列表、约束与长期价值。本章把这些输出放进可检验的比较：预测与选择质量检查输出是否符合指定标签和判断，决策效果识别检查执行策略会取得什么结果。两者都可以使用日志或实验数据，区别在于要回答的问题；时间隔离、反馈成熟、分群和资源则共同限定比较范围。读者需要据此判断，一项改进结论究竟已经获得什么证据，还缺少哪一步。

\section{评价目标与比较单位}
\label{sec:rs-evaluation-target}

评价之前应先定义\term{目标量}（estimand），即希望由数据确定的总体性质。例如，“当前用户中，点击概率预测的平均平方误差”与“向当前用户部署新列表策略后，人均有效消费增加多少”都是目标量，但它们需要的比较不同。前者固定标签和预测，后者涉及改变行动之后的结果。一个统计量是用有限数据计算的答案，不能用它的名字代替目标量。

设一次请求的可见信息为$\bm{x}$，行动为$a$，结果为$Y$。响应模型$p_\theta(\bm{x},a)$的风险可以写为
\begin{equation}
  R(\theta)
  =\mathbb E_{(\bm X,A,Y)\sim P_{\mathrm{eval}}}
    [\ell(Y,p_\theta(\bm X,A))].
  \label{eq:rs-evaluation-risk}
\end{equation}
$P_{\mathrm{eval}}$规定评价人群、行动和标签的分布。若它来自旧策略的曝光，风险就首先描述旧策略覆盖范围中的预测误差。换成新策略后，行动分布可能变化，同一个误差值不能无条件沿用。

选择质量固定的是判断集合与返回结果。给定请求的资格集合$I_t$、已判断相关对象集合$R_t$与返回列表$L_t$，命中、排序及列表覆盖分别考察哪些相关对象被找到、出现在哪些位置、是否涵盖所需方面。判断集合只覆盖一部分目录时，应同时说明未判断对象怎样处理。把未判断对象暂记为零可以定义一种基准协议，却不能由此证明它们真的不相关。

决策效果固定的则是干预。若$\pi(a\mid\bm x)$为目标策略，$Y(a)$表示执行行动$a$后的潜在结果，一步策略价值为
\begin{equation}
  V(\pi)=\mathbb E_{\bm X}
    \left[\sum_{a\in A(\bm X)}
      \pi(a\mid\bm X)\mathbb E[Y(a)\mid\bm X]\right].
  \label{eq:rs-evaluation-policy-value}
\end{equation}
$A(\bm X)$是上下文中允许的行动集合，本章此处用有限行动便于说明。比较$\pi$与$\pi_0$的增量是$\Delta=V(\pi)-V(\pi_0)$。结果可以是点击、固定窗口内购买或明确单位的收益；若行动影响未来状态，则目标应改为第~\ref{chap:rs-long-term}~章的累计回报，不能只把一步结果的名字换成“长期价值”。

\begin{table}[htbp]
  \centering
  \small
  \begin{tabularx}{\linewidth}{@{}p{22mm}XX@{}}
    \toprule
    评价任务 & 判断对象 & 证据需要说明什么\\
    \midrule
    预测质量 & 指定标签上的预测误差 & 标签含义、成熟窗口、样本分布与拟合隔离\\
    选择质量 & 指定判断集合上的候选与列表 & 资格、判断覆盖、粒度、采样和返回接口\\
    决策效果 & 策略价值或策略间增量 & 日志覆盖与识别条件，或随机分配、对照与干扰\\
    \bottomrule
  \end{tabularx}
  \caption{评价任务与证据条件。预测和选择属于质量评价；决策效果可由有条件的观测识别或随机干预检验，同一份数据可能参与多种任务。}
  \label{tab:rs-evaluation-targets}
\end{table}

表~\ref{tab:rs-evaluation-targets}中的结果都不同于转化信用归因。一次已发生的成交，按末次点击分给某个渠道，是对既成结果的记账；它没有提供“不经该渠道时会不会成交”的对照。信用规则和增量识别的具体联系留到第~\ref{chap:rs-advertising}~章，本章不把归因收入当作策略效果的替代标签。

目标量还必须匹配\term{比较单位}（unit of comparison）。对象级误差给每条对象记录一份权重，请求级指标先在每次请求内计算，用户级指标先聚合同一用户，周期级目标则可能包含多次行动。假设甲用户有100次请求、平均质量0.9，乙用户有1次请求、质量0.1。按请求平均为$90.1/101\approx0.8921$，按用户等权平均为0.5；两者没有计算矛盾，而是在回答不同人群权重的问题。把101条记录都当成101个独立用户，还会低估同一用户内相关造成的不确定性。

日志的一行、干预单位和分析单位不必相同。用户级分组可以产生很多请求记录；广告主级预算处理可以同时影响大量曝光。应保留这些对应关系，按设计确定聚合和方差估计。对同一批请求比较两个候选模型，可以计算每请求的配对质量差；若比较两套真实展示策略，则还需知道用户实际接受了哪套行动，以及行动是否改变了另一个组的机会。

一次可复核的比较因而应同时保存请求时间、可见历史、资格与目录版本、模型或策略版本、实际返回与执行结果，以及标签的事件和窗口。只保存最终均值会让后续分析无法区分模型错误、候选遗漏、资格变化与反馈未成熟。

\section{预测与选择质量评价}
\label{sec:rs-evaluation-quality}

固定目标之后，质量评价可以检查模型输出与判断之间的差异。预测质量关心单次结果及其概率，选择质量关心候选和列表；二者往往相互支持，但一个分数在某种意义下准确，不保证它适合另一种用途。尤其是对未执行行动没有标签时，质量统计只能覆盖实际获得判断的部分。

\subsection{预测质量评价}
\label{sec:rs-evaluation-prediction}

响应预测需要区分排序尺度与概率尺度。AUC衡量随机抽取的一个正例是否得分高于一个负例，并按约定给并列分数一半信用；严格单调变换不会改变这个次序。交叉熵和Brier分数直接使用预测概率，错误地报出极高置信度会改变损失。损失与概率分布差异的基础见\BookSectionRef{sec:models-probability-objectives}{概率分布的差异}，本节把它们放回同一组推荐反馈。

\begin{example}[次序相同而概率尺度不同]
\label{ex:rs-evaluation-probabilities}
教学假设有两组各100次独立曝光，点击标签均已成熟。高响应组80次点击，低响应组20次点击。模型甲对两组分别预测0.8、0.2，模型乙分别预测0.99、0.01；每组内分数相同，两模型的排序完全一致。

100个正例与100个负例共有10000对。严格正确的对为$80\times80=6400$，并列对为$80\times20+20\times80=3200$，所以两者AUC均为0.8。甲的平均交叉熵为$-0.8\log0.8-0.2\log0.2\approx0.5004$，乙为$-0.8\log0.99-0.2\log0.01\approx0.9291$。两者平均Brier分数分别为0.16和0.1961。
\end{example}

例~\ref{ex:rs-evaluation-probabilities}中，乙没有丢失次序，却把两组差异夸大了。若每次点击价值固定为10元，甲对高响应曝光的价值预测为8元，乙为9.9元；这会改变预算或出价，即使AUC一分未变。对按概率变换价值的应用，只检验区分能力是不够的。

校准检查将预测相近的样本放进同一组，比较平均预测与成熟标签频率。给定事先确定的分桶$B_1,\ldots,B_G$，常用汇总为
\begin{equation}
  \widehat{\operatorname{ECE}}
    =\sum_{g=1}^{G}\frac{|B_g|}{N}
       \left|\overline p_g-\overline y_g\right|.
  \label{eq:rs-evaluation-ece}
\end{equation}
空桶不进入求和，桶边界和最小样本量要报告。若例中的两组分别成桶，甲为0、乙为0.19。这里的0只是这份有限数据和这套分桶的结果，不是所有上下文中概率都正确的证明。更细分桶可能发现条件误差，也会增加估计噪声；只选使结果好看的分桶会引入事后选择。

概率检验还需要说明分母。曝光点击率以有效曝光为样本，点击后转化率以点击为样本，用户购买率以指定窗口中的用户为样本。标签为窗口内是否转化时，窗口结束前的零只是当前尚未观察到转化，不能与成熟负例直接混合。曝光后丢失回传的样本也需要保留缺失标记；将缺失统一填零，会把记录机制误差算成模型误差。

训练、校准和评价数据应按任务隔离。校准器在同一批标签上拟合再报告ECE，会乐观地评价自己的拟合结果；测试集反复选模型后也不再是独立验收。若目标是新对象泛化，应另外隔离对象，若目标是已有用户的下一次请求，则应保留可用历史但排除未来事件。两种划分回答不同问题，没有一套随机切分能同时代表所有部署条件。

报告总体损失后，还应检查误差分布和分群。流量大的热门广告可以主导总交叉熵，使新广告的严重高估被平均掉；以全部用户为分母和只以有点击用户为分母，也会改变解释。第~\ref{sec:rs-evaluation-time}~节统一时间与成熟条件，第~\ref{sec:rs-evaluation-groups}~节处理分群与外推；部署策略改变后会取得什么结果，仍需转到决策效果识别。

\subsection{选择质量评价}
\label{sec:rs-evaluation-selection}

选择质量把模型输出放回实际可返回的对象集合。候选命中只检查相关对象是否进入比较范围，排序质量进一步检查位置，列表覆盖还需要方面、类别或来源的定义。第~\ref{sec:rs-retrieval-evaluation}~节已建立候选诊断，本节补上使模型成绩能够相互比较的集合与判断协议。

设已判断相关对象为$R_t$，合法返回的前$k$项为$L_{t,1:k}$。召回率的分母是$|R_t|$，不是目录大小；命中率只检查交集是否非空。多个相关对象时，找到其中一个不等于全部召回。对没有已判断相关对象的请求，应预先约定不纳入该指标、单独报告覆盖，或采用另一个定义，不能在算完结果后任意决定分母。

分级相关性还要固定增益与位置折损。若第$p$位的相关性等级为$g_p$，常见的DCG按位置累计$(2^{g_p}-1)/\log_2(p+1)$，NDCG再除以同一判断集合及截断下的理想值；理想值为0时需要另行约定并报告处理规则。判断粒度从商品变为商品款式、从文档变为段落，会同时改变重复和理想排序；拆出更多片段不能被当成无代价的覆盖提高。

判断覆盖与候选范围是两个问题。即使可以给全目录打分，普通日志仍可能只知道少量对象的反馈。KuaiRec增加固定小矩阵中的消费反馈覆盖，AgentRecBench则控制工具可见信息；两者都能支持特定选择任务的比较，但数据设计不能被误写为另一种评分模型。

\subsubsection{全目录与资格集合评价协议}
\label{sec:rs-evaluation-full-catalog}

\term{全目录评价}（full-catalog evaluation）应以请求时允许访问的完整资格集合为边界。一个物品存在于离线文件，并不表示当时已经上架、仍有库存或满足地域要求。评价输入需要时间截止点、资格$I_t$、模型版本、判断与返回预算；输出保留合法候选、排名、非法对象和访问成本。

对能逐对象评分的模型，可以给$I_t$内每个对象计算分数，再按统一平局规则截取。对近似检索模型，应执行真实索引接口，不能直接用精确全量点积的理想结果替代；对标识生成模型，应执行受约束解码、标识恢复和资格检查。只接受外部候选的重排器，应同时报告该候选集的覆盖及集内排序，不把它标成全目录检索。

\begin{example}[过滤、召回与排序分开计数]
目录含背包甲、乙、丙、丁、戊、己。教学假设当前乙超过用户价格上限，己已下架，资格集合为$I_t=\{\text{甲},\text{丙},\text{丁},\text{戊}\}$。四项分数依次为0.8、0.6、0.9、0.7；当前基准只把甲标为相关，其余项按基准约定标为0。

完整资格集合次序为丁、甲、戊、丙，甲排名2，Hit@1为0、Hit@3为1，NDCG@3为$1/\log_2 3\approx0.6309$。若候选接口只返回丁、戊、丙，甲在召回阶段已经遗漏，后面的重排器无论怎样排序都无法恢复这个正例。

再假设非法对象乙、己的原始分数分别为1.1、1.0。先取整个文件前三项再过滤，只剩丁；先在资格集合内取前三项，则得到丁、甲、戊。这个差异来自过滤与截取协议，不能只归因于模型表示能力。
\end{example}

去重同样需要具体规则。一个商品的多个标题版本是否视作同一对象，引用片段是否必须排除来源文档重复，已消费物品是否还允许推荐，都应与实际任务一致。历史中出现过的物品不一定永远非法；把“未见物品推荐”的排除规则搬到复购任务会直接改变目标。

逐项评分的代价约为$O(|I_t|C_s)$，其中$C_s$为每对象评分成本；近似访问、树搜索或生成解码有自己的代价变量。全目录协议规定比较范围，不强迫所有方法采用同一计算方式。评价应记录访问、重排、过滤、补取及无结果率，使质量与实际服务能力能够对应。

\Needspace{110mm}
\subsubsection{采样指标与校正协议}
\label{sec:rs-evaluation-sampled}

大目录上逐请求评价较贵时，可以抽取部分对象参与比较。但\emph{On Sampled Metrics for Item Recommendation}表明，直接在小集合上计算命中或NDCG，可能改变模型间的相对优劣；增加重复抽样次数只减小随机波动，不能消除这种系统差异\scite{rs-e02,rs-e02-extended}

为看清原因，暂考虑一个已判断正例、$N-1$个按协议作负例的对象，固定模型分数且无并列。正例在全集合的排名为$r$。从其余对象中均匀、有放回抽取$m$次，每次抽中比正例分数更高的对象的概率为$(r-1)/(N-1)$。把重复抽到的对象按抽样次数计入比较，采样排名满足
\begin{equation}
  \widetilde r-1\mid r
    \sim\operatorname{Binomial}
       \left(m,\frac{r-1}{N-1}\right).
  \label{eq:rs-evaluation-sampled-rank}
\end{equation}
若去重或无放回采样，分布就不同；非均匀采样还依赖各对象被抽中的概率，不能沿用这个二项模型。

教学假设$N=101$、真实排名$r=12$、抽取$m=10$个对照，其中恰有一个分数更高，于是$\widetilde r=2$。全目录Hit@10为0，采样Hit@10为1；按单正例定义，全目录NDCG@10为0，采样值约0.6309。这些“负例”只是本协议的判断，不意味着未观察物品都真实无关。

对全排名函数$M(r)$，直接采样成绩的期望为
\begin{equation}
  \mathbb E[M(\widetilde r)\mid r]
    =\sum_{j=1}^{m+1}P_{r,j}M(j),
  \qquad P_{r,j}=\mathbb P(\widetilde r=j\mid r).
  \label{eq:rs-evaluation-sampled-expectation}
\end{equation}
它一般不等于$M(r)$。例中，采样Hit@10只有在十次都抽到更高对象时才为0，期望为$1-0.11^{10}$，几乎为1。样本内成绩看似稳定，仍可能几乎总在回答错误的全目录问题。

一种修正先估计真实排名：
\begin{equation}
  \widehat r
    =1+\frac{(N-1)(\widetilde r-1)}{m}.
  \label{eq:rs-evaluation-rank-correction}
\end{equation}
由二项均值可得$\mathbb E[\widehat r\mid r]=r$。例中的修正排名为11，单次不必等于12；将它代入Hit或NDCG，也不因排名无偏就获得指标无偏。非整数排名还需要预先规定取整或插值，原文实验采用向下取整。

另一种修正直接学习“每种采样排名应记多少分”。令$\bm g\in\mathbb R^{m+1}$，观测到$\widetilde r=j$就返回$g_j$。作者用所有可能真实排名上的平方偏差与方差共同确定它：
\begin{equation}
  \begin{aligned}
    \min_{\bm g}\quad&
      \sum_{r=1}^{N}
      \left[\left(\sum_jP_{r,j}g_j-M(r)\right)^2
        +\gamma v_r(\bm g)\right],\\
    v_r(\bm g)&=\sum_jP_{r,j}g_j^2-
       \left(\sum_jP_{r,j}g_j\right)^2,\qquad\gamma>0.
  \end{aligned}
  \label{eq:rs-evaluation-metric-correction}
\end{equation}
其中$v_r(\bm g)$是给定真实排名$r$时的方差，来自$\mathbb E[g_{\widetilde r}^2\mid r]-(\mathbb E[g_{\widetilde r}\mid r])^2$。这是由已知采样分布建立的二次最小化，不用测试正例的真实排名逐条训练一个回归器。形成$\bm P=[P_{r,j}]$后可解对应线性系统；矩阵奇异时需要适当的广义逆或数值正则，并披露它对结果的影响。$\gamma$控制偏差与方差取舍，应在独立验证中选择或做敏感性分析。

该计算需要约$O(Nm)$存储概率表；直接形成二次项和稠密求解还可能需要$O(Nm^2+m^3)$计算，具体可用结构或数值方法优化。固定$N,m,M,\gamma$后，评价每条记录只需查一次$g_{\widetilde r}$。但资格集合大小或采样分布变了，就不能无条件复用旧表。校正值也未必位于原指标的$[0,1]$范围，事后截断又会改变偏差。

训练负采样服务于学习目标，评价采样服务于指标估计，二者应分别记录。更好的采样校正不能补出真实相关性判断，也不保证有限数据中的所有模型排序都恢复全目录结果。计算允许时，用完整资格集合复核关键结论，才能区分采样效应与模型差异。

\Needspace{95mm}
\subsubsection{KuaiRec：近全观测反馈的选择评价}
\label{sec:rs-evaluation-kuairec}

目录可以访问而反馈不足时，评价仍会受缺失机制影响。\term{KuaiRec}在固定的小型用户--视频矩阵中增加消费反馈覆盖，使研究者能够比较稀疏观察与较完整观察下的推荐质量\scite{rs-e06}其小矩阵包含1411名用户、3327个视频，报告的交互覆盖约为99.6%；仍有用户屏蔽作者等造成的缺口，因此“近全观测”更准确地描述实际反馈范围。

数据通过约15天的补充曝光收集，而不是在同一瞬间读取所有潜在偏好。观看时长与视频时长之比可表达消费强度，原论文把观看比值超过2作为一种二元正例规则。这个阈值是评价口径，不证明多看一次必然等于喜欢，也不证明未达到阈值就没有价值。视频时长、重复观看和同一用户的历史曝光都可能影响反馈。

使用时先确定小矩阵的用户与对象范围、标签规则及缺口处理，再构造所需观察密度。若遵循论文的大矩阵训练、小矩阵评价协议，要剔除大矩阵中与小矩阵重合的交互；否则评价反馈可能提前进入训练。比较缺失机制时，可以在较完整的判断上保留不同子集作为训练观察，同时对各方法使用同一测试判断和资格范围。

教学假设某用户对甲、乙、丙三个视频的完整二元判断为$(1,1,0)$，稀疏日志只观察到甲为1、丙为0。若把乙的缺失填零，返回$(\text{乙},\text{丙})$的模型Precision@2被记为0；在较完整判断中实际是$1/2$。这说明覆盖改变了已能核验的答案。它并不说明对所有未观察对象都应反向填1。

近全观测数据还可以构造固定反馈的模拟会话，检查策略在指定状态更新规则下的选择。但补曝光过程有时间顺序，参与用户和对象又经过特定筛选，不能把小矩阵当作整个平台所有行动的潜在结果表。它补足消费反馈覆盖，对完整长期策略或市场变化的因果识别仍需另行设计。

\Needspace{110mm}
\subsubsection{AgentRecBench：工具交互与兴趣变化的基准评价}
\label{sec:rs-evaluation-agentrecbench}

推荐智能体可能先查用户记录，再查对象或评论，静态特征表无法展示它怎样取得信息。\term{AgentRecBench}把Amazon、Yelp和Goodreads中的用户、对象与评论组织成文本工具环境，评价智能体在受限可见信息下的候选选择\scite{rs-e11}它提供的是数据、工具和任务协议，模型仍需自行决定查询内容与最终排序。

任务按场景和具体请求两层过滤可见记录，并设置常规、不同历史窗口及用户或物品冷启动条件。三个月与一周历史窗口改变智能体能够看见的兴趣证据；冷启动再限制可见交互。工具返回必须遵循同一过滤，不能只从主输入删掉测试记录，却通过用户详情、对象评论或相邻记录把它取回来。

原基准的候选评价使用一个历史正例与19个未观察交互的对照对象，报告HR@1、HR@3和HR@5等命中。这里的“负例”是构造候选集时的约定，未交互不等于真实不喜欢；从20个候选中命中也不等同全目录检索。要迁移到实际服务，需要另外建立完整目录、合法输出和真实交互条件。

沿通勤例子作教学假设：三个月窗口中有六次户外装备记录，一周窗口只有两次通勤相关记录。智能体甲先读全部可见历史，再查两款背包属性；乙只读取最近事件，再查电脑夹层和价格。固定同一个目标及20项候选，两者都把目标排在第二位，HR@1为0、HR@3为1；这只能说明此例候选次序相同。若甲用了8次工具调用、乙用了3次，信息成本不同，但不能由单个例子断言乙普遍更好。

比较时应控制工具权限、模型版本、提示、调用预算与失败处理，记录调用轨迹、检索到的信息、最终候选以及时延或令牌成本。失败调用和越权返回不应悄悄从分母中删除；超过预算后如何停止，也会影响质量。工具日志使信息使用可检查，却不是策略作用于真实用户状态的实验。

历史窗口变化反映可用证据的变化，基准中的“兴趣变化”任务不能自动证明推荐行动改变了兴趣。若要研究行动如何导致未来偏好或供给响应，应回到第~\ref{chap:rs-long-term}~章的状态、策略与反馈设计；较高的离线命中率不能代替长期效果证据。

\section{决策效果识别}
\label{sec:rs-evaluation-effects}

相同的标签可以检验评分，也可以在附加条件下支持策略效果估计。区别在于：评分评价比较已经给出的预测与判断，效果识别要恢复目标行动分布下的结果。日志记录的行动受到旧策略选择，随机实验则主动控制分配；两种证据形成方式都需要明确目标策略、结果窗口与比较人群。

\subsection{观测日志下的效果识别}
\label{sec:rs-evaluation-logged}

设日志记录为$(\bm x_n,a_n,y_n,\mu_n)$，其中$\mu_n=\mu(a_n\mid\bm x_n)$是旧策略在当时信息下实际选择行动$a_n$的概率。模型打出的点击率不是这个概率，物品被用户查看的概率也不是它。若行动是完整有序列表，$\mu_n$必须描述该列表的选择；逐物品边际曝光率通常不足以重建联合概率。

在一步任务中，由日志识别式~\eqref{eq:rs-evaluation-policy-value}至少需要几项相互配合的条件。实际执行与定义的行动一致，观测结果才对应那个行动的潜在结果。给定记录的处理前信息后，行动选择还不能依赖未记录的结果原因；用潜在结果记号表示为$Y(a)\perp A\mid\bm X$。目标策略会执行的行动必须在日志中有正概率，即$\pi(a\mid\bm x)>0$时$\mu(a\mid\bm x)>0$。此外，评价人群和结果形成机制要与目标应用对应，反馈缺失或窗口截断不能被忽略。

这些条件分别处理行动含义、混杂、覆盖与迁移，不能由一个“无偏日志”标签替代。一般因果识别条件见\BookSectionRef{sec:causal-effect-identification}{因果效应的识别条件}。若旧策略把库存不足的对象排除，而新策略也必须遵守库存，就应先统一行动资格；若新策略要推荐旧日志从未出现的对象，仅靠给已有记录加权无法恢复其真实结果。

下面的推导先固定目标策略，并假设结果模型或倾向模型在独立数据上训练；需要复用数据时，可以分折拟合，在每个评价折上只使用其他折训练的辅助模型。按用户或时间具有依赖的数据，应按对应单位划分，不能让同一用户的相邻记录跨折泄漏。用评价数据同时选取收益最高的策略，还会产生选择乐观性，需要独立验收或适当的统一控制。

\Needspace{90mm}
\subsubsection{IPS：倾向加权估计}
\label{sec:rs-evaluation-ips}

旧策略偏好展示热门对象，日志自然会过多记录它们的反馈。\term{逆倾向加权}（inverse propensity weighting）用观测机制的倒数调整每条记录所代表的份额。\emph{Recommendations as Treatments}以用户--物品反馈矩阵说明这种修正怎样作用于学习与评价\scite{rs-e01}但矩阵标签风险与行动策略价值的分母和概率不同，需要分别推导。

固定有限用户集合$U$和物品集合$I$，令$O_{u,i}$表示该格反馈是否被观察，$p_{u,i}=\mathbb P(O_{u,i}=1)>0$为其纳入概率，$\ell_{u,i}$为固定预测对应的损失。完整矩阵风险及其IPS估计为
\begin{equation}
  R=\frac{1}{|U||I|}\sum_{u,i}\ell_{u,i},
  \qquad
  \widehat R_{\mathrm{IPS}}
   =\frac{1}{|U||I|}
      \sum_{u,i}\frac{O_{u,i}\ell_{u,i}}{p_{u,i}}.
  \label{eq:rs-evaluation-matrix-ips}
\end{equation}
这里的完整格数是分母，不能换成观察条数而仍称同一估计量。对固定损失取观测设计期望，$\mathbb E[O_{u,i}/p_{u,i}]=1$，所以得到$R$。这个均值等式只需要正确边际纳入概率；各格之间的依赖仍会影响方差，不能由均值无偏推出独立。

若模型也由同一份被选择的数据拟合，损失不再是上述固定量；训练风险的乐观性需要另外处理。矩阵中的纳入概率也不是点击概率，更不是完整列表的行动概率。第~\ref{sec:rs-position-bias}~节校正位置作用于点击学习的机制，与本节评价改变行动后的收益，不能共享一个未定义的“倾向”。

在上下文行动日志中，权重改为
\begin{equation}
  w_n=\frac{\pi(a_n\mid\bm x_n)}
             {\mu(a_n\mid\bm x_n)},\qquad
  \widehat V_{\mathrm{IPS}}(\pi)
       =\frac{1}{N}\sum_{n=1}^{N}w_ny_n.
  \label{eq:rs-evaluation-policy-ips}
\end{equation}
在前述条件下，给定$\bm X=\bm x$，对旧策略的行动求和有
\[
  \begin{aligned}
    \mathbb E_\mu[wY\mid\bm x]
      &=\sum_{a:\,\mu(a\mid\bm x)>0}\mu(a\mid\bm x)
         \frac{\pi(a\mid\bm x)}{\mu(a\mid\bm x)}
         \mathbb E[Y\mid\bm x,a]\\
      &=\sum_a\pi(a\mid\bm x)\mathbb E[Y(a)\mid\bm x].
  \end{aligned}
\]
概率比把行动分布换成目标分布，条件可交换性再把观测条件均值连接到潜在结果。任意删除其中一步的条件，代数消去都不足以证明真实策略价值已被恢复。

教学假设允许行动为展示甲或乙，旧策略概率为$(0.25,0.75)$，目标策略为$(0.5,0.5)$。看到一次甲被点击，权重为2，加权贡献为2；它超过二元点击的最大值1并不矛盾，贡献代表目标分布中的份额。看到乙未点击，权重为$2/3$、贡献为0。估计目标仍是总体期望，有限样本的IPS不保证落在$[0,1]$。

若四条教学日志依次为$(\text{甲},1)$、$(\text{甲},0)$、$(\text{乙},1)$、$(\text{乙},0)$，则IPS为$(2+0+2/3+0)/4=2/3$。一种自归一化形式称为\term{SNIPS}（self-normalized inverse propensity scoring）：
\begin{equation}
  \widehat V_{\mathrm{SNIPS}}
     =\frac{\sum_nw_ny_n}{\sum_nw_n}.
  \label{eq:rs-evaluation-snips}
\end{equation}
本例分母为$16/3$，结果为$1/2$。SNIPS在有界结果下能保持结果范围，但随机分子除以随机分母一般带来有限样本偏差；分母为0时无法计算，应报告覆盖失败。它也不能仅凭“归一化”保证比IPS更准确。

小概率会使极少量记录主导估计。常用诊断量为
\begin{equation}
  N_{\mathrm{eff}}
    =\frac{(\sum_n w_n)^2}{\sum_nw_n^2}.
  \label{eq:rs-evaluation-weight-ess}
\end{equation}
四条日志中它为3.2。这个量描述权重不均匀造成的信息集中，不是存在依赖、漂移或模型拟合时的通用真实样本数，更不能替代标准误。还应报告最大权重、尾部分位数、无覆盖的目标概率质量和不同裁剪阈值下的结果。

若令$\bar w=\min(w,c)$，在$Y\geq0$且其他条件正确时，$\mathbb E[\bar wY]-V(\pi)=-\mathbb E[(w-c)_+Y]$。降低大权重会引入这一向下偏差；对于可正可负的收益，偏差方向不再固定。倾向估计误差、确定性日志和未记录混杂更不能靠裁剪修好。

给定概率与结果，IPS和SNIPS都只需$O(N)$汇总，模型训练、概率重建和数据检查可能占主要成本。连续出价要用合适的条件密度比；列表或长轨迹的联合比值还可能迅速放大方差。逐请求一步加权并不会自动模拟目标策略改变预算、用户历史或市场竞争后的状态分布。

\Needspace{90mm}
\subsubsection{DR：预测与加权残差估计}
\label{sec:rs-evaluation-dr}

直接预测每个行动的结果能利用更多结构，却依赖结果模型正确；IPS减少这类建模依赖，但稀少行动的残差会被大权重放大。\term{双重稳健估计}（doubly robust estimation）将预测与加权残差结合：先在目标策略下平均预测，再用日志中实际行动的预测误差作修正\scite{rs-h13}

设$m(\bm x,a)=\mathbb E[Y\mid\bm X=\bm x,A=a]$，$\widehat m$为独立拟合的结果模型，$\widehat\mu$为日志倾向估计，参与除法的估计倾向必须为正。离散随机目标策略下的估计为
\begin{equation}
  \begin{aligned}
    \widehat V_{\mathrm{DR}}
      &=\frac1N\sum_{n=1}^{N}\psi_n,\\
    \psi_n
      &=\sum_a\pi(a\mid\bm x_n)\widehat m(\bm x_n,a)\\
      &\quad+
       \frac{\pi(a_n\mid\bm x_n)}{\widehat\mu(a_n\mid\bm x_n)}
       \left[y_n-\widehat m(\bm x_n,a_n)\right].
  \end{aligned}
  \label{eq:rs-evaluation-dr}
\end{equation}
原研究主要以确定性策略说明机制，这里把预测项写成有限随机策略的加权和。该推广的关键仍是同一个行动选择概率与同一个结果目标，不是把任意两个预测器相加。

\begin{figure}[htbp]
  \centering
  % 共享输入沿同一干线分叉；两个分量求和，不表示两套独立模型。
\begin{tikzpicture}[
  x=1mm,y=1mm,font=\figurefont\small,
  box/.style={draw=BookInk,fill=BookPaper,line width=0.8pt,
    align=center,inner sep=1.5mm,minimum height=14mm},
  flow/.style={-{Latex[length=2.2mm,width=1.5mm]},
    draw=BookInk,line width=1pt,rounded corners=1.2mm},
  wire/.style={draw=BookInk,line width=1pt,rounded corners=1.2mm}
]
  \path[use as bounding box] (0,-3) rectangle (112,73);
  \node[box,text width=94mm,fill=BookBlue!5] (input) at (56,63)
    {上下文、日志行动与结果\\目标策略、日志倾向、同一个结果模型};
  \node[box,text width=45mm] (prediction) at (26,34)
    {目标策略下的预测\\$\sum_a\pi(a\mid\bm x)\widehat m(\bm x,a)$};
  \node[box,text width=45mm] (residual) at (86,34)
    {实际行动的加权残差\\
     $\dfrac{\pi(a\mid\bm x)}{\widehat\mu(a\mid\bm x)}
       [y-\widehat m(\bm x,a)]$};
  \node[box,text width=78mm,fill=BookTeal!6] (output) at (56,6)
    {两项相加得到$\psi_n$，再按评价单位求平均\\
      $\widehat V_{\mathrm{DR}}=N^{-1}\sum_n\psi_n$};
  \draw[wire] (input.south) -- (56,49);
  \draw[flow] (56,49) -| (prediction.north);
  \draw[flow] (56,49) -| (residual.north);
  \draw[wire] (prediction.south) -- (26,20) -- (56,20);
  \draw[wire] (residual.south) -- (86,20) -- (56,20);
  \draw[flow] (56,20) -- (output.north);
\end{tikzpicture}

  \caption{双重稳健估计的两项贡献。预测项使用目标策略下所有允许行动；残差项只使用日志实际行动的结果，两个分支共享同一个结果模型。}
  \label{fig:rs-evaluation-dr}
\end{figure}

图~\ref{fig:rs-evaluation-dr}中的相加不是给模型得分任意加一个补偿。固定辅助模型并对旧策略求条件期望，估计贡献与目标条件价值之差为
\begin{equation}
  \sum_a\pi(a\mid\bm x)
    \left(1-\frac{\mu(a\mid\bm x)}
                     {\widehat\mu(a\mid\bm x)}\right)
    \left[\widehat m(\bm x,a)-m(\bm x,a)\right].
  \label{eq:rs-evaluation-dr-bias}
\end{equation}
当倾向模型正确时，第一个误差因子为0；当结果模型正确时，第二个为0。其余识别条件、支持和可积性成立时，这解释了两条一致性途径；估计模型趋近真值时还需要适当收敛条件。两者都错时，误差乘积通常不为零。“双重稳健”不表示任意两种偏差都能相互抵消。

沿前面的两行动例子，假设$\widehat m(\bm x,\text{甲})=0.6$、$\widehat m(\bm x,\text{乙})=0.3$，目标预测项为$0.5\times0.6+0.5\times0.3=0.45$。甲被点击时贡献为$0.45+2(1-0.6)=1.25$；甲未点击时为$0.45+2(0-0.6)=-0.75$。乙被点击与未点击时，贡献分别为$0.45+(2/3)(1-0.3)\approx0.9167$和$0.45-(2/3)0.3=0.25$。四条平均为$5/12\approx0.4167$。

同一批教学日志上的直接预测为0.45、IPS为0.6667、SNIPS为0.5，DR为0.4167。未提供真实总体价值时，不能根据哪个结果更接近直觉选择赢家。负的单条贡献也不是负点击，而是预测被负残差修正后的估计分量。

当结果模型接近条件均值时，残差波动通常较小；若它在稀少行动上严重失准，大权重仍会放大误差。有限样本下DR不保证优于IPS。把同一评价记录的标签拟合得残差恰为0，还可能掩盖过拟合，因此辅助模型隔离是计算协议的一部分。

有限行动数为$K$时，直接计算预测项约需$O(NK C_m)$，其中$C_m$为一次结果预测成本；残差汇总为$O(N)$。大行动集合可以用结构化求和或另行抽样近似，但会引入新的误差与随机性。支持之外的行动即使有模型预测，也主要依赖外推，不能由DR名称获得额外识别证据。

独立同分布的评价单位、固定策略及合适矩条件下，可以由$\psi_n$的样本方差构造均值标准误；同用户多次请求则需按用户聚类，长轨迹还需保留时间依赖。若用于比较两个策略，优先在同一记录上计算两套贡献之差，保留配对协方差，而不是把两个误差条当成独立再相减。

\subsection{随机干预下的效果检验}
\label{sec:rs-evaluation-randomized}

当日志覆盖或行动选择难以辨认时，可以在明确范围内随机分配策略。随机化使处理分配不依赖用户潜在结果，但只保护实际执行的分配设计；它不会自动消除处理组对对照组的影响，也不能补回未记录的结果。

设用户$u$被随机分到$Z_u\in\{0,1\}$，两组执行事先固定的策略，并在相同窗口记录$Y_u$。在稳定处理含义、完整结果记录及适当无干扰条件下，等概率分组的均值差
\begin{equation}
  \widehat\Delta
    =\frac1{n_1}\sum_{u:Z_u=1}Y_u
      -\frac1{n_0}\sum_{u:Z_u=0}Y_u
  \label{eq:rs-evaluation-experiment}
\end{equation}
估计分配这两套策略的平均效果。样本未按分配接受处理时，按分配组分析仍对应分配政策的效果；直接改为“真正点击或真正曝光者之间比较”，会按处理后的变量筛选，失去原随机对照的保护。

教学假设两组各10000名独立用户，固定窗口中分别有1200人与1000人完成有效消费。差值为0.02，即2个百分点；相对对照提升为20\%。若用二项均值的独立大样本近似，标准误为
\[
  \sqrt{\frac{0.12(1-0.12)}{10000}
             +\frac{0.10(1-0.10)}{10000}}
    \approx0.00442,
\]
95\%近似区间为$[0.0113,0.0287]$。该计算依赖用户独立、固定分析方案及充分样本；若10000条其实来自100名用户的重复请求，就不能使用这个分母宣称如此精确。

一个展示策略的实验应固定资格、流量入口与标签窗口，记录计划行动和实际执行。推荐数量、频控、延迟、失败回退与模型版本都可能属于处理的一部分。例如，新模型让服务超时增加，删掉超时请求再比较点击率，会掩盖这项完整策略的代价。结果应按预定可用性或失败口径进入统计。

随机化单位还应匹配干扰范围。同一用户跨设备进入不同组会混合处理；共享预算会使处理组多赢的机会改变对照组可买流量；提供者可能根据整个市场曝光修改供给。按用户、广告主、区域或时间段分组各有不同取舍，需要说明跨组影响、携带效应和独立单位数量。较大的集群可能减少干扰，却也减少可用于估计方差的独立单元。

实验结束条件和多指标选择应预先定义。连续查看直到出现小$p$值，再把当日当作固定终点，会改变检验误差；探索很多分群后只报告最好的，也不再是原单一检验。相关统计基础见\BookSectionRef{sec:statistics-selection-multiplicity}{选择、多重比较与持续查看}，实验设计见\BookSectionRef{sec:statistics-design-information}{对照、预算与信息获取}。未检测到差异时，还要报告区间能否排除业务上重要的损失或收益。

长观察与随机化是不同坐标。等待留存或供给响应成熟可以改变结果窗口，但不能替代对照和干扰条件；同样，长期策略也可能在充分条件下用观测日志评价。广告中的Ghost Ads与预算重分配由第~\ref{chap:rs-advertising}~章展开，本节的分配与分析原则继续适用。

\Needspace{105mm}
\subsubsection{KuaiRand：随机曝光日志的效果评价}
\label{sec:rs-evaluation-kuairand}

近乎完整地收集消费反馈与随机决定展示什么，是两个不同的数据性质。\term{KuaiRand}在真实推荐过程的一部分机会中随机注入视频，并连接用户原有行为历史与多种响应，为减少自然曝光选择偏差提供数据\scite{rs-e07}

其注入从包含7583个视频的固定池选择。一次推荐列表以某个固定概率触发随机替换，再随机选择一个槽位并均匀选择池内视频。公开材料没有给出这个触发概率的具体数值；因此，已知“条件于注入后的物品均匀”不能直接写出整个系统执行某完整列表的概率。

条件于一次有效注入及其固定槽位，若视频均匀选取，每个池内视频的选择概率为$1/7583$。但用户进入、列表生成、槽位选择、注入触发及历史筛选是更完整的机制。要用式~\eqref{eq:rs-evaluation-policy-ips}评价另一策略，必须先判断目标行动是否正好是这个条件实验中的视频选择；不能把$1/7583$当成任意对象在任意请求中的曝光倾向。

公开数据还按至少有10次随机曝光筛选用户，保留27285名用户。这个条件由一段观察期内的记录决定，外推到所有用户时需要注意纳入机制；随机注入后的局部对象分配与公开用户样本具有代表性，是两项不同要求。数据中的12种反馈也不是同一种标签，点击、观看及其他交互应各自定义窗口和分母。

用两类视频说明选择效应。教学假设自然曝光中，热门视频100次有30次正响应，冷门视频20次有2次；条件随机注入中，两类各50次，分别有10次和9次正响应。自然记录的30\%与10\%，同时混入了用户、位置和选择；随机机会中的20\%与18\%是在该条件分配下更可比的结果，但有限样本差异仍有不确定性。这些构造数字不来自KuaiRand实验，也不能外推成“冷门视频应全面增加曝光”。

应用时应区分正常日志与注入日志，核验可见历史、随机标记、视频池、槽位及反馈。训练可以使用允许的正常历史，评价则保留注入机制的条件；若用注入反馈调参，还需要独立验收子集。用户序列中的前一次注入可能影响后续行为，分析多步结果时不能把所有随机曝光都当成无历史影响的独立试验。

KuaiRand增加了特定池和机会中的行动覆盖，可支持相应局部响应或策略比较。它不能自动识别全目录、任意列表、完整长期策略或供给生态的效果。与KuaiRec配合时，应分别说明使用近全反馈解决什么缺失问题、使用随机分配解决什么比较问题，而不是将两者统称为无限范围的“无偏数据”。

\section{共同评价条件}
\label{sec:rs-evaluation-conditions}

前面的估计各自固定了目标和证据条件，真实比较还会受到时间、人群与资源变化影响。更长的历史、更多的调用或不同的用户构成，可能比评分结构本身更强地改变成绩。共同条件需要随结果保存，才能判断改善来自哪项变化。可靠性、隐私和数据撤回又对方案提出额外要求，不能由平均命中率统一代替。

\subsection{时间与反馈成熟度}
\label{sec:rs-evaluation-time}

一次交互至少有事件发生时间、系统获知时间和标签达到观察窗口终点的时间。历史$H_{u,t}$应只包含请求时已经能够取得的内容；某条事件发生在请求之前，但直到请求之后才回传，也不能提前作为特征。对象属性、聚合统计、特征字典和嵌入更新都需要遵守这一可用时间，而不只是原始点击行。

时间隔离与标签成熟处理不同问题。时间隔离阻止预测使用未来信息，成熟窗口规定何时能够给结果作判断。假设10月1日发生点击，转化窗口为7天；10月3日尚未转化时，只知道前两天没有转化。把它立即记成七日负例，会把“尚未知道”误成“最终没有”；10月6日发生的转化也不能在10月3日训练时提前加入。

训练截止与评价截止应同时给出。评价未来请求时，模型只能使用当时已发布的数据；汇总其七日结果时，需要再等待七日并考虑回传延迟。若只纳入截止时已经转化的样本，会偏向短延迟正例；若将未成熟者全标负，会偏向低估。可以等待统一成熟、明确只估计短窗口结果，或在有依据的删失与延迟模型下估计更长窗口，三者应按不同协议报告。

时间漂移还会使早期人群与后期人群不同。模型版本、目录规模和活动期间都可能改变误差。滚动评价应给出每个时期的结果、样本数和可用信息，不能只保留表现最好的时间段。观察更久能够增加结果信息，却不能自动修复未观察混杂、行动无覆盖或市场干扰。

\Needspace{95mm}
\subsubsection{全局时间线评价协议}
\label{sec:rs-evaluation-timeline}

逐用户把最后一次交互留作测试，看似每人都只用自己的过去，却仍可能读取别人的未来。\term{全局时间线评价}（global timeline evaluation）要求在所有用户共享的时钟上检查可用信息。相关研究指出，按用户各自切分会让未来交互和未来对象进入更早请求的训练或候选目录\scite{rs-e03}

图~\ref{fig:rs-evaluation-timeline}给出一个教学反例：用户甲的最后交互在3日，用户乙的最后交互在5日。逐用户最后留出时，乙在4日的行为被分入训练，却已经晚于甲在3日的测试。若丙物品到5日才上架，在整个数据文件中提前建立目录也会泄漏未来供给。

\begin{figure*}[htbp]
  \centering
  % 教学日序；方框标记交互，虚线是3日请求之前的信息边界。
\begin{tikzpicture}[
  x=1mm,y=1mm,font=\figurefont\small,
  event/.style={draw=BookInk,fill=BookBlue!5,line width=0.8pt,
    text width=21mm,minimum height=11mm,align=center,inner sep=1mm},
  future/.style={event,draw=BookAmber,dashed,fill=BookAmber!6},
  time/.style={-{Latex[length=2.2mm,width=1.5mm]},
    draw=BookMuted,line width=0.8pt},
  note/.style={font=\figurefont\footnotesize,text=BookMuted}
]
  \path[use as bounding box] (0,-2) rectangle (169,67);
  \foreach \x/\day in {30/1,60/2,90/3,120/4,150/5}{
    \node at (\x,53) {\day 日};
  }
  \node[anchor=west] at (0,35) {用户甲};
  \node[anchor=west] at (0,14) {用户乙};
  \draw[time] (20,35) -- (167,35);
  \draw[time] (20,14) -- (167,14);
  \draw[draw=BookTeal,dashed,line width=1pt] (76,0) -- (76,59);
  \node[anchor=south,text=BookTeal] at (76,61) {3日请求前的信息边界};
  \node[event] at (30,35) {1日交互\\可用历史};
  \node[event,fill=BookTeal!8,draw=BookTeal] at (90,35)
    {3日请求\\本次测试};
  \node[event] at (60,14) {2日交互\\可用历史};
  \node[future] at (120,14) {4日交互\\此时不可用};
  \node[future] at (150,14) {5日交互\\乙的测试};
  \node[note,anchor=north] at (135,5) {物品丙也到5日才上架};
\end{tikzpicture}

  \caption{两用户交错时间的教学示意。以3日甲的请求为评价时点，2日及之前已获知的记录可以使用，4日乙的交互和5日才上架的对象均不可使用。虚线只标记信息截止，事件仍按各自发生和回传时间判定。}
  \label{fig:rs-evaluation-timeline}
\end{figure*}

部署式全局截止选择一个共同$t_0$，用此前已可见数据训练固定模型，再评价之后的请求。滚动更新则在每个指定时点，用此前已释放的标签更新模型，并评价下一个区间。它们都能遵守全局可用性，但更新频率与陈旧程度不同，不能只写“按时间划分”便认为协议等价。

论文的时间线方案将全部记录按时间组织，在其中标记测试点，沿时间推进；每个测试点预测时只使用允许的既有记录。若逐事件都先预测、后在可用时释放结果给学习器，就得到一种极细粒度的\term{预序贯评价}（prequential evaluation）。它要求算法支持相应更新；对必须批量重训的模型，可以采用明确的窗口近似，但要报告额外滞后及维护成本。

实施时先排序请求与数据发布事件，再维护当前历史、目录资格和模型状态。到请求事件时保存模型版本并预测；到标签发布事件时，才按协议允许训练。点击可能很快释放，七日转化或月留存则需排入未来可用队列。测试标签在成熟后用于后续训练并不自动违规，关键是评价目标是否明确允许在线适应，以及后续请求是否确实发生在标签已知之后。

沿图中的例子，3日请求可用甲1日、乙2日且已回传的记录。4日收到乙的新交互后，才可将它用于以后请求；5日丙上架后，才可进入资格集合。若甲3日交互的转化要到10日才成熟，5日模型仍不能使用其完整转化标签。按事件日期排序而忽略发布时间，也可能留下隐蔽泄漏。

全局排序通常需要$O(N\log N)$预处理，已按时间记录时可以流式合并。此后的预测、更新、索引重建与成熟队列成本取决于方法，不能只报告最终离线打分时延。协议输出应包含每个请求的可用历史范围、目录和模型版本、预测及标签释放时点，便于复查。

时间合规只证明没有使用指定时间边界之外的信息。一个完全按时间执行的旧策略仍可能只覆盖热门物品；一个长期运行的实验仍可能有跨组干扰。这些问题分别回到第~\ref{sec:rs-evaluation-logged}~节和第~\ref{sec:rs-evaluation-randomized}~节处理。

\subsection{分群与应用范围}
\label{sec:rs-evaluation-groups}

总体均值需要知道由谁贡献。新用户、低频用户、新对象、不同来源和发生漂移的时期，可能同时具有不同反馈覆盖与不同模型误差。分群应在比较前按可用的处理前信息定义；用实验后的活跃度分组，可能把策略造成的活跃变化当成固定人群差异。

\begin{table}[htbp]
  \centering
  \small
  \begin{tabularx}{\linewidth}{@{}Xrrr@{}}
    \toprule
    用户组 & 用户权重 & 原方案质量 & 新方案质量\\
    \midrule
    高频用户 & 0.20 & 0.80 & 0.90\\
    低频用户 & 0.80 & 0.50 & 0.49\\
    等权用户总体 & 1.00 & 0.560 & 0.572\\
    \bottomrule
  \end{tabularx}
  \caption{总体改善与分群下降可以同时发生的教学数据。组别按比较前活跃度固定，组内指标与用户权重口径一致；表中不含抽样不确定性。}
  \label{tab:rs-evaluation-groups}
\end{table}

表~\ref{tab:rs-evaluation-groups}中总体提高0.012，低频组却下降0.01，两组差距由0.30扩大为0.41。是否可以接受，取决于第~\ref{chap:rs-constraints}~章已确定的多方目标和底线，而不是总体均值自动决定。若改按请求加权，高频用户的权重又会变化，必须重新说明目标人群。

每组至少应报告规模、有效反馈覆盖、质量或效应差和不确定性。少数新对象的均值波动较大，不能仅凭一个极端点估计宣布模型失效；反过来，样本不足也不能当作它们没有受到损害的证明。若许多组都被检查，应区分预先指定的核心组和事后探索，并处理多重比较。

组内相关性与分母同样重要。一个商家的大量相似商品不等于同样数量的独立提供者；同一用户的重复曝光也不等于独立用户。可采用匹配设计的聚类标准误或重采样，但必须保留抽样与干预结构。区间和分位数的含义见\BookSectionRef{sec:statistics-intervals-quantiles}{区间、分位数与预测范围}。

目标人群改变时，如果组内结果机制可迁移且各组有覆盖，可以按新的组权重重新汇总；若新用户群在旧数据中完全缺失，重加权无法生成该群的结果。跨域搜索、新查询与不同广告预算阶段还需相应任务条件，分别由后两章承接。分群使边界可见，不会单独带来因果识别。

\subsection{信息与计算资源}
\label{sec:rs-evaluation-resources}

两种方法的输入相同，才有机会把差异归于内部计算。一个方法使用近期文字、图像和外部候选，另一个仅使用用户ID，即使后者更慢或更差，也不能只据此比较某个网络层。模型结构、可用信息和服务接口应分别记录；实际系统比较允许它们一起变化，但结论应对应整套方案。

目录检索、候选重排和直接生成尤其容易混比。目录检索要支付索引维护与候选访问成本，重排器依赖外部候选覆盖，生成器需要解码、标识恢复、去重和资格核验。只记录神经网络前向时间，会漏掉不同接口中最重要的代价。PRM、NAR4Rec等展示方法的输入边界见第~\ref{sec:rs-prm}~节与第~\ref{sec:rs-nar4rec}~节。

教学假设三种方案在同一目录和硬件上得到表~\ref{tab:rs-evaluation-resources}。这些数字只说明记账方式：表中的重排器若不计候选生成，会显得比完整服务便宜；生成器若不计核验，也会遗漏非法标识的成本。

\begin{table}[htbp]
  \centering
  \small
  \begin{tabularx}{\linewidth}{@{}Xrrr@{}}
    \toprule
    方案 & 候选或解码 & 评分与核验 & 请求合计\\
    \midrule
    目录检索后排序 & 4\,ms & 6\,ms & 10\,ms\\
    外部候选上重排 & 4\,ms & 3\,ms & 7\,ms\\
    标识生成后核验 & 8\,ms & 4\,ms & 12\,ms\\
    \bottomrule
  \end{tabularx}
  \caption{服务阶段计时的教学口径。各行包含相同的外部等待范围，数字不是实测性能；只有在候选覆盖、输出质量与失败处理同时满足要求时，合计时延才可用于方案比较。}
  \label{tab:rs-evaluation-resources}
\end{table}

平均时延之外，还应检查尾部、并发负载、缓存冷热与失败率。训练时间、索引存储、定期更新和撤回后重建属于周期成本，不能随意摊入或移出单请求指标。智能体工具调用还包括远端等待、令牌和重复查询，用户能得到的实际结果取决于预算耗尽后的停止规则。

结构消融应尽量固定信息、候选和调优预算；完整系统比较则应报告实际可执行方案的质量与总成本。如果新方法用更多信息或更高预算取得更高质量，可以如实描述这种交换，但不能把全部收益归于一个层的数学形式。更多计算也不能补足未覆盖行动的真实结果，资源条件始终不替代证据协议。

\subsection{可靠性、隐私与数据撤回评价}
\label{sec:rs-evaluation-protection}

平均预测质量之外，推荐系统还需判断何时证据不足、泄漏了什么信息，以及指定数据撤回后原训练影响是否仍在。可靠性、隐私和撤回各自有对象和判据，可以同时约束同一模型，却不能合并成一个没有单位的“可信分数”。

可靠性评价需要检查预测范围和不确定性是否支持实际决策。一个模型拒绝20\%请求后，在其余80\%上有90\%正确率，只能说明已接受部分的质量；必须同时报告覆盖和拒绝人群。用不确定性阈值筛选后，还要在未参与阈值选择的数据上验收。第~\ref{sec:rs-distribution-free-reliability}~节的集合风险保证依赖校准分布和完整选择规则，重新增加多样性后处理或迁移人群可能改变保证。

概率与范围检验见\BookSectionRef{sec:reliability-output-evaluation}{概率与预测范围的检验}，拒答与风险筛选见\BookSectionRef{sec:reliability-selection-evaluation}{风险识别与拒答效果的检验}。分桶校准、区间覆盖和集合错误控制检验不同性质，任意一个通过都不能替代其他对象的检查。

隐私评价需要固定攻击者能访问什么：模型参数、嵌入、梯度、接口概率或只有推荐列表。成员推断、属性推断和重建攻击各自对应不同泄漏。经验攻击失败只表明指定攻击未成功，不证明对所有攻击安全。使用差分隐私时，应核验相邻数据关系、裁剪、噪声和组合后的预算，再同时报告效用与成本；基础定义见\BookChapterRef{chap:trust-privacy}{基础、理论和可信性卷的“隐私保护”章}。

撤回的目标不是从日志里删掉一行，而是处理该记录已经进入的表示、参数、统计与索引。评价要明确删除单位是交互、用户还是物品，比较对象是用保留数据重新训练的模型分布，还是另一种近似标准。保留数据上的质量、对被撤回对象的影响和完整维护时延应一起检查；恶意污染造成的异常又是另一问题，相关方法见第~\ref{sec:rs-robust-scoring}~节与第~\ref{sec:rs-unlearning-maintenance}~节。

\Needspace{100mm}
\subsubsection{CURE4Rec：推荐撤回的多维检验}
\label{sec:rs-evaluation-cure4rec}

只比较撤回后的命中率，可能让一个保留了大量原训练信息的模型显得很好。\term{CURE4Rec}主要面向用户级推荐撤回，同时检查撤回完整性、保留任务效用、活跃度与分片之间的公平性，以及运行成本\scite{rs-e12}它是一套评价基准，不是把任何模型变成可撤回模型的通用算法。

协议先确定原训练集合、撤回用户和保留集合，再分别得到近似撤回结果与相应参考。重新训练参考需要沿用模型族、调参口径与随机性处理，不能因为一次随机初始化不同就把所有差异归于删除失败。近似撤回还可能通过协同关系影响未删除用户，因此不能只查看被删除用户的参数是否置零。

选择撤回对象时，原工作区分核心与边缘，其相关性分数使用用户度与邻居物品度均值的乘积。教学假设用户甲连接4个物品，物品度为$(10,20,10,20)$，分数为$4\times15=60$；用户乙只连接2个物品，度为$(40,50)$，分数为$2\times45=90$。乙交互更少，仍可能有更高的该项核心分数，故不能把核心直接改写为最活跃。

完整性的一种经验检查训练成员推断器，输入用户嵌入及其交互物品的平均嵌入，判断目标用户是否仍像原训练成员。构造成员与非成员样本、拟合攻击器和最终评价应相互隔离，并保持两类用户的活跃度等条件可比。否则，攻击器可能只学会识别高活跃用户，而非残留训练影响。

若攻击AUC从0.80降到0.52，且重训参考为0.51，这只能作为该攻击下更接近参考的教学观察；没有样本量和区间，不能断言0.52与0.51等效，更不能证明严格撤回。更换攻击器、访问权限或未来查询后，残留影响可能仍被发现。

原基准使用HR与NDCG@20检查保留任务，并按保留用户中最活跃的前5\%与其余95\%比较组均值；对于分片式方案，还检查分片效用方差。这些指标分别反映质量、组间影响与分片不均，不能只挑其中改善的一项作为整体通过。删除用户可能同时改变邻居的嵌入和候选覆盖，应按完整推荐接口重新评价。

运行成本应从请求到完成验收计量。教学假设近似撤回更新需40秒，相关表示修复30秒、重建索引20秒、验证10秒，共100秒；完整重训需400秒，另加20秒索引与10秒验证，共430秒。完整流程的速度比为4.3，而不是只用400除以40得到10。若批量撤回、缓存失效或服务回退不同，还应分别报告吞吐、尾部时延和期间质量。

最终结果应同时给出删除单位、请求数量与分布、参考模型、攻击访问条件、保留质量、分群影响和完整耗时。用户级撤回通过的经验检查，不能直接外推到任意物品或单条交互删除；对严格保证的要求仍由算法定义与证明承担。

\section{本章小结}
\label{sec:rs-evaluation-summary}

评价的第一步是确定希望知道什么。预测误差检查指定标签上的数值，选择质量检查明确判断集合中的返回结果，策略价值与因果增量则比较执行行动后的后果。相同数据可以参与不同任务，但必须重新确认目标、分母、单位和识别条件，不能从一项标签成绩直接跳到部署收益。

完整资格集合与采样协议决定排名的比较范围，反馈覆盖决定哪些判断已被观察。IPS用行动概率比改变日志代表的分布，DR进一步利用结果预测与残差；它们都需要支持与识别条件。随机干预主动控制分配，却仍要核验实际执行、结果记录与跨组影响。近全反馈、局部随机曝光和工具环境各自补充一部分证据，没有一种资源提供无限范围的效果保证。

时间线与反馈成熟使数据符合请求时的知识边界，分群说明均值背后的得失，资源条件说明方法能否在相同约束下执行。可靠性、隐私与撤回则检验额外要求。一项方案值得采用，需要这些证据与它实际改变的行动相匹配，而不是仅获得某一个更高的数字。

有了共同评价语言，后两章可以把判断落到更具体的应用。搜索需要核验结果是否满足当前明确需求，广告还要识别资金、支付和竞争改变后的效果；它们将沿用本章的目标量、可见信息与比较协议，并补充各自特有的条件。
